\section{\texorpdfstring{Example: \(2\times 2\) matrices --- a genuinely noncommutative ring}{Example: 2\textbackslash times 2 matrices --- a genuinely noncommutative ring}}\label{sec:08-rings:07-matrices:1}

\texttt{intRing} above is commutative (\texttt{Int.mul\_\allowbreak{}comm} exists, even though it was not wired into \texttt{Ring}'s fields), so it does not test the fact that \texttt{Ring} deliberately does \emph{not} assume \texttt{mul\_\allowbreak{}comm}. Matrices do: \(M_2(\mathbb{Z})\), the ring of \(2\times 2\) integer matrices under matrix addition and multiplication, has \(AB \neq BA\) in general. It is also a good test of \href{https://loogle.lean-lang.org/?q=mul_assoc}{\texttt{mul\_\allowbreak{}assoc}}, since matrix multiplication's associativity is not a one-line library lemma call. It genuinely takes some work, which is the point of working through it.

A \(2 \times 2\) matrix is represented as a \texttt{structure} with four entries. (More scalable representations exist --- \texttt{Fin 2 → Fin 2 → Int}, or Mathlib's general \href{https://loogle.lean-lang.org/?q=Matrix}{\texttt{Matrix}} --- but four named fields keep every computation fully explicit, which is the goal here.)

\begin{lstlisting}
structure Mat2 where
  a11 : Int
  a12 : Int
  a21 : Int
  a22 : Int

-- Core Lean 4 does not auto-generate a field-wise extensionality lemma
-- for a plain `structure` (that convenience is a Mathlib `@[ext]`
-- attribute); we supply one by hand, using the `mk.injEq` lemma core Lean
-- *does* generate for every structure, so it can be used below exactly
-- like an auto-generated `.ext` would be.
theorem Mat2.ext {X Y : Mat2} (h1 : X.a11 = Y.a11) (h2 : X.a12 = Y.a12)
    (h3 : X.a21 = Y.a21) (h4 : X.a22 = Y.a22) : X = Y := by
  cases X
  cases Y
  rw [Mat2.mk.injEq]
  exact ⟨h1, h2, h3, h4⟩
\end{lstlisting}

\begin{mathreading}

\texttt{Mat2} is the free \(\mathbb{Z}\)-module \(M_2(\mathbb{Z}) \cong \mathbb{Z}^4\) on the four matrix entries. The extensionality lemma \texttt{Mat2.ext} --- supplied by hand right alongside the structure, since almost every proof below needs it --- says two \texttt{Mat2} values are equal exactly when all four entries match.

\end{mathreading}

\begin{lstlisting}
def Mat2.add (X Y : Mat2) : Mat2 where
  a11 := X.a11 + Y.a11
  a12 := X.a12 + Y.a12
  a21 := X.a21 + Y.a21
  a22 := X.a22 + Y.a22
\end{lstlisting}

\texttt{Mat2.add} is entrywise \(+\).

\begin{lstlisting}
def Mat2.neg (X : Mat2) : Mat2 where
  a11 := -X.a11
  a12 := -X.a12
  a21 := -X.a21
  a22 := -X.a22
\end{lstlisting}

\texttt{Mat2.neg} is entrywise \(-\).

\begin{lstlisting}
def Mat2.zero : Mat2 := ⟨0, 0, 0, 0⟩
\end{lstlisting}

\texttt{Mat2.zero} is the zero matrix.

\begin{lstlisting}
-- (row i, col k) entry of X * Y is Σⱼ X[i,j] * Y[j,k]; with only two
-- indices this sum is just two terms, written out directly.
def Mat2.mul (X Y : Mat2) : Mat2 where
  a11 := X.a11 * Y.a11 + X.a12 * Y.a21
  a12 := X.a11 * Y.a12 + X.a12 * Y.a22
  a21 := X.a21 * Y.a11 + X.a22 * Y.a21
  a22 := X.a21 * Y.a12 + X.a22 * Y.a22
\end{lstlisting}

\texttt{Mat2.mul} is the matrix product with \((XY)_{ik} = \sum_j X_{ij}Y_{jk}\), here the two-term sums since \(n = 2\).

\begin{lstlisting}
def Mat2.one : Mat2 := ⟨1, 0, 0, 1⟩
\end{lstlisting}

\texttt{Mat2.one} is the identity \(I = \begin{psmallmatrix}1&0\\0&1\end{psmallmatrix}\). Together these five definitions are the operations of the matrix ring \(M_2(\mathbb{Z})\).

\subsection{Why matrix multiplication is noncommutative --- check it computationally first}\label{sec:08-rings:07-matrices:2}

Before building the \texttt{Ring Mat2} instance, consider the very fact that motivates this example:

\begin{lstlisting}
def X : Mat2 := ⟨1, 1, 0, 1⟩
def Y : Mat2 := ⟨1, 0, 1, 1⟩

#eval Mat2.mul X Y   -- ⟨2, 1, 1, 1⟩
#eval Mat2.mul Y X    -- ⟨1, 1, 1, 2⟩
\end{lstlisting}

\texttt{\#eval} here is doing real work: it is a two-line proof by computation that \texttt{mul} is not commutative. This is cheaper than any hand-written counterexample proof, and it is exactly the kind of thing to try before committing to a \texttt{theorem}. Should \texttt{¬ ∀ X Y, Mat2.mul X Y = Mat2.mul Y X} be needed as a theorem, this computed pair \texttt{(X, Y)} serves as the witness.

\begin{mathreading}

This is the standard counterexample to commutativity of matrix multiplication: with \(X =
\begin{psmallmatrix}1&1\\0&1\end{psmallmatrix}\) and \(Y =
\begin{psmallmatrix}1&0\\1&1\end{psmallmatrix}\) (elementary shear matrices), \[
XY = \begin{psmallmatrix}2&1\\1&1\end{psmallmatrix} \neq
\begin{psmallmatrix}1&1\\1&2\end{psmallmatrix} = YX,
\] so \(M_2(\mathbb{Z})\) is a \emph{non}commutative ring. The pair \((X, Y)\) is a concrete witness for the existential \(\exists X, Y,\ XY \neq YX\).

\end{mathreading}

\textbf{Mathlib equivalent.} Mathlib's \texttt{Matrix (Fin 2) (Fin 2) Int} is exactly \(M_2(\mathbb{Z})\), already a (noncommutative) \texttt{Ring} instance. There is no \texttt{Mat2}/ \texttt{Mat2.ext}/\texttt{add4\_\allowbreak{}reorder} needed, and \href{https://lean-lang.org/doc/reference/latest/Tactic-Proofs/Tactic-Reference/}{\texttt{decide}} works directly here (unlike \texttt{Mat2}) because \texttt{Matrix} over \texttt{Int} already has \texttt{DecidableEq}:

\begin{lstlisting}
example : Ring (Matrix (Fin 2) (Fin 2) Int) := inferInstance

def X' : Matrix (Fin 2) (Fin 2) Int := !![1, 1; 0, 1]
def Y' : Matrix (Fin 2) (Fin 2) Int := !![1, 0; 1, 1]

example : X' * Y' ≠ Y' * X' := by decide
\end{lstlisting}

\subsection{The additive group and the ring}\label{sec:08-rings:07-matrices:3}

\begin{lstlisting}
def mat2Group : Group Mat2 where
  op := Mat2.add
  id := Mat2.zero
  inv := Mat2.neg
  assoc := by
    intro X Y Z
    -- Goal: Mat2.add (Mat2.add X Y) Z = Mat2.add X (Mat2.add Y Z)
    show Mat2.mk _ _ _ _ = Mat2.mk _ _ _ _
    -- Each of the four field-equations reduces to Int addition's
    -- associativity, entrywise. `Mat2.ext` (the extensionality lemma
    -- defined above: two Mat2's are equal iff all four fields match) lets
    -- us split the single Mat2 equality into four Int equalities.
    apply Mat2.ext <;> exact Int.add_assoc _ _ _
  id_left := by
    intro X
    apply Mat2.ext <;> exact Int.zero_add _
  id_right := by
    intro X
    apply Mat2.ext <;> exact Int.add_zero _
  inv_left := by
    intro X
    apply Mat2.ext <;> exact Int.add_left_neg _
  inv_right := by
    intro X
    apply Mat2.ext <;> exact Int.add_right_neg _
\end{lstlisting}

\texttt{mat2Group} verifies that \((M_2(\mathbb{Z}), +)\) is a group: every axiom (\texttt{assoc}, \texttt{id\_\allowbreak{}left}, \texttt{id\_\allowbreak{}right}, \texttt{inv\_\allowbreak{}left}, \texttt{inv\_\allowbreak{}right}) is checked entrywise, using \texttt{Mat2.ext} to split the single \texttt{Mat2} equality into four \texttt{Int} equalities, each of which is then a direct citation of the matching \texttt{Int} addition lemma.

\begin{lstlisting}
def mat2CommGroup : CommGroup Mat2 where
  toGroup := mat2Group
  comm := by
    intro X Y
    apply Mat2.ext <;> exact Int.add_comm _ _
\end{lstlisting}

\texttt{mat2CommGroup} upgrades \texttt{mat2Group} to an abelian group by supplying commutativity, again entrywise via \texttt{Mat2.ext} and \texttt{Int.add\_\allowbreak{}comm}.

\begin{lstlisting}
-- A reusable shuffle: rearranges (a+b)+(c+d) into (a+c)+(b+d) — exactly
-- what's needed whenever a "product of sums" expands into four cross
-- terms that must be regrouped to match the other side.
theorem add4_reorder (a b c d : Int) : a + b + (c + d) = a + c + (b + d) := by
  rw [Int.add_assoc a b (c + d)]
  rw [show b + (c + d) = c + (b + d) from by
    rw [← Int.add_assoc, Int.add_comm b c, Int.add_assoc]]
  rw [← Int.add_assoc a c (b + d)]

def mat2Ring : Ring Mat2 where
  addGrp := mat2CommGroup
  mul := Mat2.mul
  one := Mat2.one
  mul_assoc := by
    intro X Y Z
    -- Unlike addition, each entry of (X * Y) * Z expands to a sum of
    -- *four* products of three matrix entries (vs. Int.mul_assoc's single
    -- rebracketing) — associativity of matrix multiplication is
    -- distributivity-plus-associativity of the underlying ring applied
    -- four times per entry, not a single lemma call. Each entry equation
    -- unfolds (via `Int.add_mul`/`Int.mul_add`/`Int.mul_assoc`) to two
    -- sums of the same four products in a different order; `add4_reorder`
    -- is exactly the regrouping needed to match them up.
    apply Mat2.ext
    · show (X.a11 * Y.a11 + X.a12 * Y.a21) * Z.a11 + (X.a11 * Y.a12 + X.a12 * Y.a22) * Z.a21
          = X.a11 * (Y.a11 * Z.a11 + Y.a12 * Z.a21) + X.a12 * (Y.a21 * Z.a11 + Y.a22 * Z.a21)
      rw [Int.add_mul, Int.add_mul, Int.mul_add, Int.mul_add,
          Int.mul_assoc, Int.mul_assoc, Int.mul_assoc, Int.mul_assoc,
          add4_reorder]
    · show (X.a11 * Y.a11 + X.a12 * Y.a21) * Z.a12 + (X.a11 * Y.a12 + X.a12 * Y.a22) * Z.a22
          = X.a11 * (Y.a11 * Z.a12 + Y.a12 * Z.a22) + X.a12 * (Y.a21 * Z.a12 + Y.a22 * Z.a22)
      rw [Int.add_mul, Int.add_mul, Int.mul_add, Int.mul_add,
          Int.mul_assoc, Int.mul_assoc, Int.mul_assoc, Int.mul_assoc,
          add4_reorder]
    · show (X.a21 * Y.a11 + X.a22 * Y.a21) * Z.a11 + (X.a21 * Y.a12 + X.a22 * Y.a22) * Z.a21
          = X.a21 * (Y.a11 * Z.a11 + Y.a12 * Z.a21) + X.a22 * (Y.a21 * Z.a11 + Y.a22 * Z.a21)
      rw [Int.add_mul, Int.add_mul, Int.mul_add, Int.mul_add,
          Int.mul_assoc, Int.mul_assoc, Int.mul_assoc, Int.mul_assoc,
          add4_reorder]
    · show (X.a21 * Y.a11 + X.a22 * Y.a21) * Z.a12 + (X.a21 * Y.a12 + X.a22 * Y.a22) * Z.a22
          = X.a21 * (Y.a11 * Z.a12 + Y.a12 * Z.a22) + X.a22 * (Y.a21 * Z.a12 + Y.a22 * Z.a22)
      rw [Int.add_mul, Int.add_mul, Int.mul_add, Int.mul_add,
          Int.mul_assoc, Int.mul_assoc, Int.mul_assoc, Int.mul_assoc,
          add4_reorder]
  one_mul := by
    intro X
    apply Mat2.ext
    · show 1 * X.a11 + 0 * X.a21 = X.a11
      rw [Int.one_mul, Int.zero_mul, Int.add_zero]
    · show 1 * X.a12 + 0 * X.a22 = X.a12
      rw [Int.one_mul, Int.zero_mul, Int.add_zero]
    · show 0 * X.a11 + 1 * X.a21 = X.a21
      rw [Int.zero_mul, Int.one_mul, Int.zero_add]
    · show 0 * X.a12 + 1 * X.a22 = X.a22
      rw [Int.zero_mul, Int.one_mul, Int.zero_add]
  mul_one := by
    intro X
    apply Mat2.ext
    · show X.a11 * 1 + X.a12 * 0 = X.a11
      rw [Int.mul_one, Int.mul_zero, Int.add_zero]
    · show X.a11 * 0 + X.a12 * 1 = X.a12
      rw [Int.mul_zero, Int.mul_one, Int.zero_add]
    · show X.a21 * 1 + X.a22 * 0 = X.a21
      rw [Int.mul_one, Int.mul_zero, Int.add_zero]
    · show X.a21 * 0 + X.a22 * 1 = X.a22
      rw [Int.mul_zero, Int.mul_one, Int.zero_add]
  left_distrib := by
    intro X Y Z
    apply Mat2.ext
    · show X.a11 * (Y.a11 + Z.a11) + X.a12 * (Y.a21 + Z.a21)
          = (X.a11 * Y.a11 + X.a12 * Y.a21) + (X.a11 * Z.a11 + X.a12 * Z.a21)
      rw [Int.mul_add, Int.mul_add, add4_reorder]
    · show X.a11 * (Y.a12 + Z.a12) + X.a12 * (Y.a22 + Z.a22)
          = (X.a11 * Y.a12 + X.a12 * Y.a22) + (X.a11 * Z.a12 + X.a12 * Z.a22)
      rw [Int.mul_add, Int.mul_add, add4_reorder]
    · show X.a21 * (Y.a11 + Z.a11) + X.a22 * (Y.a21 + Z.a21)
          = (X.a21 * Y.a11 + X.a22 * Y.a21) + (X.a21 * Z.a11 + X.a22 * Z.a21)
      rw [Int.mul_add, Int.mul_add, add4_reorder]
    · show X.a21 * (Y.a12 + Z.a12) + X.a22 * (Y.a22 + Z.a22)
          = (X.a21 * Y.a12 + X.a22 * Y.a22) + (X.a21 * Z.a12 + X.a22 * Z.a22)
      rw [Int.mul_add, Int.mul_add, add4_reorder]
  right_distrib := by
    intro X Y Z
    apply Mat2.ext
    · show (X.a11 + Y.a11) * Z.a11 + (X.a12 + Y.a12) * Z.a21
          = (X.a11 * Z.a11 + X.a12 * Z.a21) + (Y.a11 * Z.a11 + Y.a12 * Z.a21)
      rw [Int.add_mul, Int.add_mul, add4_reorder]
    · show (X.a11 + Y.a11) * Z.a12 + (X.a12 + Y.a12) * Z.a22
          = (X.a11 * Z.a12 + X.a12 * Z.a22) + (Y.a11 * Z.a12 + Y.a12 * Z.a22)
      rw [Int.add_mul, Int.add_mul, add4_reorder]
    · show (X.a21 + Y.a21) * Z.a11 + (X.a22 + Y.a22) * Z.a21
          = (X.a21 * Z.a11 + X.a22 * Z.a21) + (Y.a21 * Z.a11 + Y.a22 * Z.a21)
      rw [Int.add_mul, Int.add_mul, add4_reorder]
    · show (X.a21 + Y.a21) * Z.a12 + (X.a22 + Y.a22) * Z.a22
          = (X.a21 * Z.a12 + X.a22 * Z.a22) + (Y.a21 * Z.a12 + Y.a22 * Z.a22)
      rw [Int.add_mul, Int.add_mul, add4_reorder]
\end{lstlisting}

Two points are worth noting:

\begin{enumerate}
\def\labelenumi{\arabic{enumi}.}
\tightlist
\item
  \textbf{Every proof obligation here is spelled out explicitly, with no automation} --- matching the rest of the book, and unlike an earlier draft of this section, which reached for Mathlib's \texttt{ring} tactic (a decision procedure for commutative-ring identities). Since this book never imports Mathlib, \texttt{ring} is not actually available. Each entry equation is instead unfolded by hand through \texttt{Int.add\_\allowbreak{}mul}/\texttt{Int.mul\_\allowbreak{}add}/\texttt{Int.mul\_\allowbreak{}assoc} down to a sum of the same four cross terms in a different order, and \texttt{add4\_\allowbreak{}reorder} (proved once, above, and reused twelve times) supplies exactly the regrouping needed to match them. The \texttt{Ring Mat2} bundle itself is still noncommutative --- nothing here decides \emph{that} automatically. \texttt{mul := Mat2.mul} is supplied directly, and the \texttt{mul\_\allowbreak{}comm}-shaped fact is simply absent from \texttt{Ring}'s fields, exactly as Chapter 8's exercise on \texttt{left\_\allowbreak{}distrib}/\texttt{right\_\allowbreak{}distrib} anticipated.
\item
  \textbf{This is the general pattern for ``ring of \(n\times n\) matrices over a commutative ring \(S\)'':} the entries live in \(S\), every \texttt{Ring Mat2} proof obligation reduces (through distributivity/associativity in \(S\), applied a bounded number of times depending on \(n\)) to a polynomial identity purely in \(S\), and \(\text{Mat}_n(S)\) itself is noncommutative as soon as \(n \geq 2\), regardless of whether \(S\) is commutative. Matrix rings are the standard first example that any general ring theory (Chapter 9's theorems, Mathlib's \texttt{Ring} hierarchy) has to handle \emph{without} assuming commutativity, which is exactly why \texttt{Ring} does not bake in \texttt{mul\_\allowbreak{}comm}.
\end{enumerate}

\begin{mathreading}

\texttt{mat2Ring} is the ring \(M_2(\mathbb{Z})\) assembled formally: \texttt{mat2Group}/\texttt{mat2CommGroup} verify that \((M_2(\mathbb{Z}), +)\) is an abelian group (all axioms hold \emph{entrywise}, which is what \texttt{Mat2.ext} followed by the \(\mathbb{Z}\)-lemma expresses --- \(M_2(\mathbb{Z}) \cong
\mathbb{Z}^4\) as additive groups), and \texttt{mat2Ring} supplies the monoid \((M_2(\mathbb{Z}), \times, I)\) with distributivity. Associativity of \(\times\) is the one substantial fact: \(((XY)Z)_{i\ell} = \sum_{j,k} X_{ij}Y_{jk}Z_{k\ell}
= (X(YZ))_{i\ell}\), whose per-entry form is a polynomial identity over the commutative ring \(\mathbb{Z}\), so \texttt{ring} can close it. The construction generalizes exactly the same way to \(M_n(S)\) for any commutative ring \(S\): entrywise reduction turns each ring axiom into an \(S\)-polynomial identity, and \(M_n(S)\) is noncommutative for \(n \ge 2\).

\end{mathreading}

\textbf{Mathlib equivalent, continued.} Where the book spends most of this section deriving \texttt{mul\_\allowbreak{}assoc} for \texttt{Mat2} by hand (the \texttt{add4\_\allowbreak{}reorder} helper, reused twelve times across all five axioms), Mathlib already proves associativity of matrix multiplication generically. It is simply \texttt{mul\_\allowbreak{}assoc} again, the same lemma name as every other ring in this chapter's Mathlib boxes:

\begin{lstlisting}
example (A B C : Matrix (Fin 2) (Fin 2) Int) : (A * B) * C = A * (B * C) :=
  mul_assoc A B C
\end{lstlisting}

And where the book's \texttt{mat2Ring} needs a custom \texttt{Int}-arithmetic rewrite chain for each of \texttt{mul\_\allowbreak{}assoc}/\texttt{one\_\allowbreak{}mul}/\texttt{mul\_\allowbreak{}one}/\texttt{left\_\allowbreak{}distrib}/ \texttt{right\_\allowbreak{}distrib}, Mathlib has an automated decision procedure for exactly this class of goal, \href{https://lean-lang.org/doc/reference/latest/Tactic-Proofs/Tactic-Reference/}{\texttt{noncomm\_\allowbreak{}ring}} --- the noncommutative-ring counterpart of the \texttt{ring} tactic that the book's own note above states is not available without Mathlib:

\begin{lstlisting}
example (A B C : Matrix (Fin 2) (Fin 2) Int) :
    A * (B + C) = A * B + A * C := by noncomm_ring
\end{lstlisting}
